La primera versión pasó siete de los ocho chequeos. Falló F6: el orden temporal observado
en el caso A, con \texttt{ORD=2}, fue \textbf{1{,}45}, con diferencias sucesivas
$1{,}47\cdot10^{-7}$ y $5{,}39\cdot10^{-8}$. Como la puerta no se edita, se buscó la causa.
La secuencia de controles, en orden:

\begin{enumerate}
\item \textbf{No es el fondo no-deslizante}: con fondo free-slip el orden es 1{,}48.
\item \textbf{No es el esquema de SPECTER en general}: los caminos \texttt{freeslip} y
  \texttt{noslip} planos, sobre un modo 3D con $w\neq0$, dan orden $\approx2{,}4$ en la
  energía.
\item \textbf{No es el acople de $\eta$}: partiendo de un modo de velocidad suave, la
  superficie deformable da orden $\approx3$ con $g=0$ y con $g=1$.
\item \textbf{No es la condición inicial incompatible}. Arrancar del reposo con
  $\eta\neq0$ viola la compatibilidad (en $t=0^+$, $\partial_tu=-\nabla p_0$ no respeta
  el no-deslizamiento ni la tensión nula) y se esperaba una reducción de orden por las
  capas límite impulsivas. Pero el mismo RK2 aplicado a la referencia MAC, con la misma
  condición inicial y el mismo $\Delta z$, da orden \textbf{2{,}00} exacto. La hipótesis
  quedó refutada.
\item \textbf{Es un error por paso.} La diferencia entre SPECTER y el RK2 de la referencia
  (extrapolada en el espacio) debería ser independiente de $dt$ si SPECTER fuera un RK2
  consistente de alguna discretización espacial. No lo es: crece al achicar $dt$
  ($8{,}07$, $8{,}29$, $8{,}61$, $9{,}20\cdot10^{-7}$ para
  $dt=10^{-3}\dots1{,}25\cdot10^{-4}$), aproximadamente en proporción al número de pasos.
  Con $N_z=64$ domina todo: las diferencias sucesivas se \emph{duplican} al dividir $dt$
  por dos.
\item \textbf{Aislado con $dt\to0$}: con $dt=10^{-9}$ y $10^{-10}$ el miembro derecho es
  nulo y el estado no debería cambiar. En los caminos \texttt{freeslip} y
  \texttt{freesurface} la energía de \archivo{balance.txt} cambia $-1{,}4\cdot10^{-7}$ por
  paso con $N_z=64$ y $\approx-10^{-8}$ con $N_z=128$, \emph{igual} para los dos $dt$; en
  \texttt{noslip} el cambio es proporcional a $dt$, el físico.
\end{enumerate}

\paragraph{La causa.} Los caminos \texttt{freeslip} y \texttt{freesurface} necesitan
$v^*_z$ en la cara, y el parche de la Fase~2 lo conseguía mandando $v_z$ por
\texttt{goto\_domain\_w\_boundaries} y de vuelta por \texttt{goto\_3d\_fourier}. Esa vuelta
\emph{recalcula la continuación FC-Gram} de $v_z$ a partir de sus valores físicos, y después
de una proyección ---que actúa sobre todo el dominio periódico extendido--- la
continuación de $v_z$ ya no coincide con la que FC-Gram le daría a sus valores físicos. Cada
subpaso cambia entonces el campo en una cantidad que \emph{no se anula con $dt$}. El camino
\texttt{noslip} nunca transforma $v_z$, y por eso no lo tiene. Un primer intento de
corrección, suponiendo que el defecto era de la reconstrucción de Neumann del valor de
borde, no cambió nada y se descartó.

\paragraph{La corrección.} $v^*_z$ se queda en el dominio $(k_z,k_y,k_x)$ y sólo se leen sus
trazas en las caras (\texttt{fs\_trace}). Con eso el cambio por paso a $dt\to0$ pasa a ser
proporcional a $dt$ ($-2{,}2\cdot10^{-9}$ y $-2{,}2\cdot10^{-10}$ con $N_z=128$), y el orden
del caso A es \textbf{2{,}00}, con diferencias $1{,}381\cdot10^{-7}$ y
$3{,}450\cdot10^{-8}$, a 6\,\% del RK2 puro de la referencia ($1{,}458\cdot10^{-7}$ y
$3{,}646\cdot10^{-8}$).

\paragraph{Lo que implica para la Fase 2\propio.} El camino \texttt{freeslip\_z} plano, que
ya se usó en las corridas del barrido en $\delta$ y del puente de banda ancha, \emph{tiene
el mismo defecto}. La puerta de la Fase~2 no podía verlo: sus chequeos de tasa (V4, V5)
usan el modo de corte sin estructura horizontal, que no lo tiene (el modo $k=0$ de $v_z$ es
nulo), y V3 mide la tensión, no la energía. Su tamaño depende de la resolución vertical y del
flujo: para el modo de prueba ($k_x=1$) es $\sim10^{-8}$ de la energía por paso con
$N_z=128$ y $\sim1{,}4\cdot10^{-7}$ con $N_z=64$, que es la malla de
\archivo{decaimiento_banda_ancha.py}. Esos números son de la energía de
\archivo{balance.txt}, que \texttt{energy()} suma sólo sobre los puntos \emph{físicos}
(\archivo{pseudospec_hd.f90}, lazos \texttt{k = 1,nz-Cz}): \textbf{es disipación espuria en
el dominio físico}. Una primera versión de este informe decía lo contrario, que era la energía
del dominio extendido; lo refutó el verificador leyendo esa rutina y se comprobó. Sesga las
tasas hacia arriba, y \etq{P-02} es justamente una comparación de tasas.

\textbf{Cuánto afecta a la Fase~4 no se midió.} Una estimación de esta sesión, sin validar:
la deriva por unidad de tiempo es $\delta/dt$, con $\delta$ el cambio por paso. Con
$\delta=1{,}4\cdot10^{-7}$ y $dt=5\cdot10^{-4}$ serían $2{,}8\cdot10^{-4}$ por unidad de
tiempo del código, contra una tasa de energía $2\alpha\approx0{,}2$ en esas unidades: un
$0{,}1$\,\% para el modo de prueba, que tiene $w\neq0$. Las corridas de la Fase~4 arrancan
con $w=0$ (flujo horizontal cuasi-2D), donde el mecanismo sólo actúa a través del $v_z$ que
genera la no linealidad, así que el sesgo ahí debería ser menor. Es una estimación, no una
medida.
\paragraph{Corregido también en el camino plano \etq{D-50}.} Lucía decidió corregir
\texttt{freeslip\_z}: recibe la traza de $v^*_z$ en la cara (\texttt{trace\_z}) y $v_z$ ya no
hace la vuelta. Antes y después, con el mismo modo de prueba:

\begin{center}\footnotesize
\begin{tabular}{lccc}
\toprule
 & \multicolumn{2}{c}{deriva por paso, $dt=10^{-9}$ / $10^{-10}$} & $N_z=64$: diferencias \\
 & $N_z=64$ & $N_z=128$ & entre $dt$ sucesivos\\
\midrule
antes & $-1{,}44$ / $-1{,}43\cdot10^{-7}$ & $-1{,}05\cdot10^{-8}$ / $-8{,}4\cdot10^{-9}$ & $2{,}5\to3{,}9\to6{,}7\cdot10^{-6}$\\
después & $-9{,}4\cdot10^{-10}$ / $-8{,}0\cdot10^{-11}$ & $-2{,}25\cdot10^{-9}$ / $-2{,}25\cdot10^{-10}$ & $9{,}6\to7{,}1\to4{,}2\cdot10^{-7}$\\
\bottomrule
\end{tabular}
\end{center}

Después de la corrección la deriva es proporcional a $dt$ (la física) y las diferencias
convergen. Lo que se temía para la Fase~4 se midió: el puente de banda ancha 16² de
\etq{V-19}, vuelto a correr con la misma configuración, cambia sus tasas en
$\le5\cdot10^{-10}$ relativo (el control lineal, en $10^{-15}$). Esas corridas arrancan con
$w=0$, y el defecto sólo actúa a través del $v_z$ en la cara. \textbf{Los resultados de la
Fase~4 quedan como están}; el 32² de \etq{V-19}, \etq{V-16} y \etq{V-17}, de la misma clase,
no se volvieron a correr.
